\section{Introduction}
\label{sec:intro} 
Healthcare professionals have faced burnout for decades (\cite{marzanoDiagnosingOvercrowdedEmergency2024,batandaPrevalenceBurnoutHealthcare2024,bolatovOccupationalBurnoutHealthcare2025,tangBurnoutStressNew2025,haldaneHealthSystemsResilience2021,arsenaultCOVID19ResilienceHealthcare2022}).
Amidst routine heavy case loads where physicians must take history, examine, investigate, and treat patients, documentation comes in as an additional stress to physicians (\cite{preiksaitisChatGPTNotSolution2023,levyDefiningDocumentationBurden2024,muradMeasuringDocumentationBurden2024}).
As new medical technologies emerge, they often increase rather than decrease workload.
For instance, computed tomography angiography began being performed directly in emergency departments (EDs) for patients with suspected acute aortic pathology.
This eliminated the need to transport patients to radiology with the goal of improving efficiency.
However, the proportion of detected aortic pathology cases remained the same.
Meanwhile, the number of images and series per scan increased significantly, adding a substantial burden to radiologists (\cite{kanzariaEmergencyPhysicianPerceptions2015,sokolovskayaEffectFasterReporting2015,pourvaziriImagingInformationOverload2022,brulsWorkloadRadiologistsOncall2020,danlantsmanTrendRadiologistWorkload2022}).
The 21st Century Cures Act (Cures Act) further added to the burden on physicians, especially radiologists, as it forces them to update reports and notes to the EHR as early as possible (\cite{mehanImmediateRadiologyReport2022,choudhryEnhancedReadabilityDischarge2019}).
In many countries, radiology workspaces are still not optimized for work efficiency, which adds to the burden (\cite{hillmanRadiologistsBurdenInefficiency2013,krupinskiLongRadiologyWorkdays2010}). 
All these factors affect the quality of life of radiologists and the workflow of physicians dependent on radio-diagnosis, hence the efficacy of treatment for patients as well (\cite{atwalHealthIssuesRadiologists2017,groverPsychologicalProblemsBurnout2018}).
The above issue of burnout concerning radiologists could be viewed as one tree in a forest, considering the vastness of the prevalent healthcare system (\cite{rotensteinPrevalenceBurnoutPhysicians2018}).
Identifying the critical pathway (documentation in this context) and working on it would ultimately lead to a better outcome (\cite{everyCriticalPathwaysReview2000}).
With the advancements in large language models (LLMs), there have been many works done with clinicians, engineers, and researchers to reduce the burden of documentation (\cite{jeblickChatGPTMakesMedicine2023,patelChatGPTFutureDischarge2023,sunshineEvaluatingQualityUnderstandability2025,thirunavukarasuLargeLanguageModels2023,singhalLargeLanguageModels2023}).
Additionally, Medical Information Mart for Intensive Care (MIMIC)-IV (\cite{johnsonMIMICIV,johnsonMIMICIVNoteDeidentifiedFreetext}), an extensive publicly available database, has been used to study the stereotypical biases exhibited by LLMs (\cite{bolukbasiManComputerProgrammer2016}).
With advancements in LLMs, there are many papers that show how these models are aiding in text generation and summarization in general (\cite{vanveenAdaptedLargeLanguage2024,chenBenchmarkingLargeLanguage2025,buschCurrentApplicationsChallenges2025}).
While multiple studies demonstrate how LLMs can help summarize or generate medical documentation (progress notes and radiology reports), a critical question remains unanswered: do clinicians actually communicate these radiology findings to patients? (\cite{siewertImpactCommunicationErrors2016,europeansocietyofradiologyesrPatientSurveyValue2021,rockallPatientCommunicationRadiology2022,kasalakWorkOverloadDiagnostic2023})
In most cases, radiology reports route through the referring physician, who then may or may not talk to the patient about all the radiology findings, as the referring physician would focus more on moving to the next level in the patient care pipeline.
There are huge chances for miscommunication or lack of communication happening in this workflow (\cite{vandenbergPatientComplaintsRadiology2019,vijanOptimizingPatientCommunication2023}).
LLMs offer significant potential to aid in bridging this communication gap between patients and radiologists by providing patient-centric explanations of radiology reports as a first-pass solution.